% GENERATED FILE. Edit canonical lessons, then run npm run generate:books.
% canonical-source-hash: fnv1a64:a3d6a63e24bd6a6b
% canonical-lessons: HI-C495-kasrat, HI-C495-yog, HI-C495-parhez, HI-C495-paustik, HI-C495-pacna

\chapter{Eating Well, Keeping Fit}
\label{ch:hi-eating-well-keeping-fit}
\hlchaptermodality{495}

\begin{chapteropening}
\textbf{By the end of this chapter:} \emph{I can say what I do to stay fit.}

Complete the last lesson of chapter 495: I can say what I do to stay fit.
\end{chapteropening}

\section[\texorpdfstring{\hi{कसरत} (kasrat)}{कसरत (kasrat)}]{\hi{कसरत} (kasrat) — exercise; a workout}

\label{lesson:HI-C495-kasrat}



\begin{quote}
Before the new one: say the Hindi for to warm (a sore part), then the Hindi for to cover someone (with a blanket).
\end{quote}

\subsection*{You'll want to know: \hi{कसरत}}
\textbf{\hi{कसरत}} — \emph{kasrat} — \textquotedblleft{}exercise; a workout\textquotedblright{}.

\begin{quote}
\hi{अरुण} \hi{रोज़} \hi{सुबह} \hi{आधा} \hi{घंटा} \hi{कसरत} \hi{करता} \hi{है।}
\end{quote}

\begin{quote}
\emph{Aruṇ roz subah ādhā ghaṇṭā kasrat kartā hai.}
\end{quote}

\begin{quote}
Arun exercises for half an hour every morning.
\end{quote}

Feminine.

\begin{cousinweb}
\textbf{\hi{कसरत}} (\emph{kasrat}) is usually traced to Arabic \emph{kathra}, \textquotedblleft{}abundance, a great number\textquotedblright{}, which reached Hindi through Persian \emph{kasrat}. The sense \textquotedblleft{}physical exercise\textquotedblright{} is the one Hindi and Urdu use; how it grew from \textquotedblleft{}abundance\textquotedblright{} is not settled.
\end{cousinweb}

\begin{grammarlens}[title={What you've built}]
One more word for describing what you lived and felt.
\end{grammarlens}

\subsection*{Your turn}
\begin{itemize}
\raggedright
  \item \emph{Say it:} \emph{kasrat}
  \item \emph{Say it:} \emph{kasrat}, once more
  \item \emph{Say it:} \emph{Aruṇ roz subah ādhā ghaṇṭā kasrat kartā hai.} — Arun exercises for half an hour every morning.
  \item \emph{Recall:} say the Hindi for to warm (a sore part), then the Hindi for to cover someone (with a blanket), then say \emph{kasrat} again
\end{itemize}

\begin{tcolorbox}[breakable,colback=teal!4,colframe=teal!35!black,title={Before you move on}]
What does \hi{कसरत} mean? (\textquotedblleft{}Exercise; a workout\textquotedblright{}.) Say it once more.
\end{tcolorbox}
\section[\texorpdfstring{\hi{योग} (yog)}{योग (yog)}]{\hi{योग} (yog) — yoga}

\label{lesson:HI-C495-yog}



\begin{quote}
Before the new one: say the Hindi for to cover someone (with a blanket), then the Hindi for exercise.
\end{quote}

\subsection*{You'll want to know: \hi{योग}}
\textbf{\hi{योग}} — \emph{yog} — \textquotedblleft{}yoga\textquotedblright{}.

\begin{quote}
\hi{दादी} \hi{रोज़} \hi{योग} \hi{करती} \hi{हैं}, \hi{इसलिए} \hi{वह} \hi{चुस्त} \hi{हैं।}
\end{quote}

\begin{quote}
\emph{dādī roz yog kartī haiṁ, isliye vah cust haiṁ.}
\end{quote}

\begin{quote}
Grandmother does yoga every day, which is why she is so spry.
\end{quote}

Masculine.

\begin{cousinweb}
\textbf{\hi{योग}} (\emph{yog}) is a tatsama, borrowed from Sanskrit \emph{yoga}, first \textquotedblleft{}the act of yoking, joining\textquotedblright{} (Monier-Williams); English took \emph{yoga} from the same word.
\end{cousinweb}

\begin{grammarlens}[title={What you've built}]
One more word for describing what you lived and felt.
\end{grammarlens}

\subsection*{Your turn}
\begin{itemize}
\raggedright
  \item \emph{Say it:} \emph{yog}
  \item \emph{Say it:} \emph{yog}, once more
  \item \emph{Say it:} \emph{dādī roz yog kartī haiṁ, isliye vah cust haiṁ.} — Grandmother does yoga every day, which is why she is so spry.
  \item \emph{Recall:} say the Hindi for to cover someone (with a blanket), then the Hindi for exercise, then say \emph{yog} again
\end{itemize}

\begin{tcolorbox}[breakable,colback=teal!4,colframe=teal!35!black,title={Before you move on}]
What does \hi{योग} mean? (\textquotedblleft{}Yoga\textquotedblright{}.) Say it once more.
\end{tcolorbox}
\section[\texorpdfstring{\hi{परहेज़} (parhez)}{परहेज़ (parhez)}]{\hi{परहेज़} (parhez) — avoiding certain foods; abstinence}

\label{lesson:HI-C495-parhez}



\begin{quote}
Before the new one: say the Hindi for exercise, then the Hindi for yoga.
\end{quote}

\subsection*{You'll want to know: \hi{परहेज़}}
\textbf{\hi{परहेज़}} — \emph{parhez} — \textquotedblleft{}avoiding certain foods; abstinence\textquotedblright{}.

\begin{quote}
\hi{मधुमेह} \hi{में} \hi{मीठे} \hi{से} \hi{परहेज़} \hi{करना} \hi{पड़ता} \hi{है।}
\end{quote}

\begin{quote}
\emph{madhumeh meṁ mīṭhe se parhez karnā paṛtā hai.}
\end{quote}

\begin{quote}
With diabetes, you have to keep off sweet things.
\end{quote}

Masculine.

\begin{cousinweb}
\textbf{\hi{परहेज़}} (\emph{parhez}) is usually traced to Persian \emph{parhēz}, \textquotedblleft{}abstinence, keeping away\textquotedblright{}. \emph{X se parhez karnā} is to keep off X.
\end{cousinweb}

\begin{grammarlens}[title={What you've built}]
One more word for describing what you lived and felt.
\end{grammarlens}

\subsection*{Your turn}
\begin{itemize}
\raggedright
  \item \emph{Say it:} \emph{parhez}
  \item \emph{Say it:} \emph{parhez}, once more
  \item \emph{Say it:} \emph{madhumeh meṁ mīṭhe se parhez karnā paṛtā hai.} — With diabetes, you have to keep off sweet things.
  \item \emph{Recall:} say the Hindi for exercise, then the Hindi for yoga, then say \emph{parhez} again
\end{itemize}

\begin{tcolorbox}[breakable,colback=teal!4,colframe=teal!35!black,title={Before you move on}]
What does \hi{परहेज़} mean? (\textquotedblleft{}Avoiding certain foods; abstinence\textquotedblright{}.) Say it once more.
\end{tcolorbox}
\section[\texorpdfstring{\hi{पौष्टिक} (pauṣṭik)}{पौष्टिक (pauṣṭik)}]{\hi{पौष्टिक} (pauṣṭik) — nourishing}

\label{lesson:HI-C495-paustik}



\begin{quote}
Before the new one: say the Hindi for yoga, then the Hindi for avoiding certain foods.
\end{quote}

\subsection*{You'll want to know: \hi{पौष्टिक}}
\textbf{\hi{पौष्टिक}} — \emph{pauṣṭik} — \textquotedblleft{}nourishing\textquotedblright{}.

\begin{quote}
\hi{दाल} \hi{और} \hi{सब्ज़ी} \hi{पौष्टिक} \hi{खाना} \hi{है।}
\end{quote}

\begin{quote}
\emph{dāl aur sabzī pauṣṭik khānā hai.}
\end{quote}

\begin{quote}
Lentils and vegetables are nourishing food.
\end{quote}

\begin{cousinweb}
\textbf{\hi{पौष्टिक}} (\emph{pauṣṭik}) is a tatsama, borrowed from Sanskrit \emph{pauṣṭika}, \textquotedblleft{}nourishing, invigorating\textquotedblright{}, from \emph{puṣṭi}, \textquotedblleft{}growth\textquotedblright{} (Monier-Williams).
\end{cousinweb}

\begin{grammarlens}[title={What you've built}]
One more word for describing what you lived and felt.
\end{grammarlens}

\subsection*{Your turn}
\begin{itemize}
\raggedright
  \item \emph{Say it:} \emph{pauṣṭik}
  \item \emph{Say it:} \emph{pauṣṭik}, once more
  \item \emph{Say it:} \emph{dāl aur sabzī pauṣṭik khānā hai.} — Lentils and vegetables are nourishing food.
  \item \emph{Recall:} say the Hindi for yoga, then the Hindi for avoiding certain foods, then say \emph{pauṣṭik} again
\end{itemize}

\begin{tcolorbox}[breakable,colback=teal!4,colframe=teal!35!black,title={Before you move on}]
What does \hi{पौष्टिक} mean? (\textquotedblleft{}Nourishing\textquotedblright{}.) Say it once more.
\end{tcolorbox}
\section[\texorpdfstring{\hi{पचना} (pacnā)}{पचना (pacnā)}]{\hi{पचना} (pacnā) — to be digested}

\label{lesson:HI-C495-pacna}



\begin{quote}
Before the new one: say the Hindi for avoiding certain foods, then the Hindi for nourishing.
\end{quote}

\subsection*{You'll want to know: \hi{पचना}}
\textbf{\hi{पचना}} — \emph{pacnā} — \textquotedblleft{}to be digested\textquotedblright{}.

\begin{quote}
\hi{भारी} \hi{खाना} \hi{जल्दी} \hi{नहीं} \hi{पचता।}
\end{quote}

\begin{quote}
\emph{bhārī khānā jaldī nahīṃ pactā.}
\end{quote}

\begin{quote}
Heavy food is slow to digest.
\end{quote}

Past (it …): \hi{वह} \hi{खाना} \hi{पच} \hi{गया}, \emph{vah khānā pac gayā}.

\begin{cousinweb}
\textbf{\hi{पचना}} (\emph{pacnā}) is a tadbhava: CDIAL lists Hindi \emph{pācnā}, \textquotedblleft{}to cook\textquotedblright{}, and from it the intransitive \emph{pacnā}, \textquotedblleft{}to be digested\textquotedblright{}, under Sanskrit \emph{pacyátē}, \textquotedblleft{}is cooked, is digested\textquotedblright{} (CDIAL 7654). \emph{pacānā}, \textquotedblleft{}to digest\textquotedblright{}, is its transitive partner.
\end{cousinweb}

\begin{grammarlens}[title={What you've built}]
One more word for describing what you lived and felt.
\end{grammarlens}

\subsection*{Your turn}
\begin{itemize}
\raggedright
  \item \emph{Say it:} \emph{pacnā}
  \item \emph{Say it:} \emph{pacnā}, once more
  \item \emph{Say it:} \emph{bhārī khānā jaldī nahīṃ pactā.} — Heavy food is slow to digest.
  \item \emph{Recall:} say the Hindi for avoiding certain foods, then the Hindi for nourishing, then say \emph{pacnā} again
\end{itemize}

\begin{tcolorbox}[breakable,colback=teal!4,colframe=teal!35!black,title={Before you move on}]
What does \hi{पचना} mean? (\textquotedblleft{}To be digested\textquotedblright{}.) Say it once more.
\end{tcolorbox}
